%% sections/abstract.tex

We study the Z2 quantum Mpemba effect (QME) — the phenomenon in which a more strongly
Z2-symmetry-broken quantum state restores its symmetry faster than a less broken companion —
in the one-dimensional transverse-field Ising model and the anisotropic XY chain.
Initialising the system in a product spin state with polarisation angle $\theta$, and
monitoring the Z2 entanglement asymmetry $\DS_A(t)$ of a subsystem following a sudden quench,
we find that Z2 QME is a universal and generic feature of integrable dynamics:
it is present across all 23 post-quench field values tested, spanning the ferromagnetic
phase (no dynamical quantum phase transitions, DQPTs), the quantum critical point, and the
paramagnetic phase (DQPTs present).
This stands in sharp contrast to chaotic random circuits, where Z2 QME is absent.
An exact closed-form expression for the initial entanglement asymmetry,
$\DS_A(0;\theta) = \hbin\!\bigl((1+\cos^{\NA}\!\theta)/2\bigr)$, anchors the analysis
and saturates the universal Z2 upper bound.
Statistical analysis of 671 initial-state pairs reveals that the crossing time $\tM$
is strongly correlated with the initial EA imbalance (Spearman $\rho = +0.90$,
$p \approx 3 \times 10^{-240}$), while the distribution of $\tM$ within each DQPT period
is non-uniform but not centred on DQPT singularities (KS $p \approx 6 \times 10^{-25}$;
binary proximity test $p = 0.27$): DQPTs modulate but do not govern crossing times.
The effect persists in the thermodynamic limit (finite-size scaling over $N = 8$--$18$)
and is robust across the XY chain family ($\gamma = 0$--$1$).
These results identify the initial entanglement asymmetry geometry as the
master variable controlling Z2 QME in integrable systems, providing both a quantitative
prediction and a direct handle for quantum simulation platforms.
